\section{引言}
热带气旋能够改变森林结构与鸟类生境，其影响往往延续到风暴结束之后。波多黎各已有研究记录了飓风造成的森林损伤和后续鸟类响应，为理解扰动相关变化在岛屿景观中的差异提供了基础\cite{uriarte2019,wunderle1995}。近期研究在考虑探测概率后发现，波多黎各平鸽在飓风后仍经历长期种群瓶颈；针对波多黎各红树林的遥感研究则显示，不同地点的冠层损伤与植被指数恢复存在差异\cite{rivera2025,howe2025}。这些研究提示，环境比较需要分别考察植被变化和鸟类响应，不能预设两者具有相同的恢复轨迹。对于岛屿特有留鸟，低报告时期之后的记录增加，需要结合其发生环境来理解。分析既要确定报告在何处回升，也要检验海拔和植被绿度是否与原有占域地点在下一时期仍被占据的概率相关，进而判断这些关系能否指导其他地点的监测。

波多黎各绿蜂鸟（\Rm）为检验这条证据链提供了研究对象。Ortiz-Andrade 和 Rojas-Perez\cite{ortiz2026}结合 eBird 清单、海拔和卫星植被绿度，分析由 María 与 Fiona 两次飓风界定的三个时期。其清单层模型将报告概率与环境梯度联系起来，动态占域模型则在不完全探测条件下考察占域持续性。公开的可复现包使我们能够重新检查这两类互补证据，也使“复现已发表数值”与“确认实现对应预期统计模型”成为可以分别回答的问题。

第一个解释难点在于响应变量的含义。eBird 清单记录的是报告事件，其概率同时受到物种存在、调查努力及其他观测条件影响\cite{sullivan2009,johnston2021}。Backstrom 等\cite{backstrom2025}进一步表明，eBird 数据中的空间与时间采样偏差能够相互作用并影响趋势估计，因此时期比较还需检查样本的环境组成。因此，阳性清单比例的变化本身并不代表种群数量变化。不同海拔的预测能够揭示全岛平均值掩盖的差异，但仍然是条件报告概率。动态占域模型则在重复调查与探测假设下，估计地点的潜在占域状态及其在占据和未占据之间的转变\cite{mackenzie2002,mackenzie2003,guillera2017}。近期草原林莺研究结合长期定点调查与动态占域模型，分别评价了定殖和局地灭绝对生境管理的响应\cite{gordon2024}，说明已占据地点的保留与新地点进入占据状态需要分别分析。两类结果各有解释范围；出现得分较高或其补数较低，不能直接作为局地灭绝转移的测量。

第二个难点在于支持潜在状态推断的观测过程是否得到正确实现。持续性的环境关联需要以观测及其协变量正确对应为前提。一个可复现的脚本也可能稳定地复现非预期排列。检查具体实现问题，有助于判断生态解释是否可靠。校正可能保留环境关联，同时改变探测差异的证据，因此需要分别检验不同解释。

第三个难点是地理预测。概括已采样地点的关系，未必能在留出的地理区块中保留相同价值。重复访问和地理聚集使验证设计必须与监测用途相匹配\cite{roberts2017,valavi2019,ploton2020}。随机划分与空间划分回答不同问题，一种方案下的较好表现不能自动推广到另一种方案。近期物种分布模型比较进一步发现，随机验证可能高估模型的推广表现\cite{spacetime2025}；空间与时间环境关系的对照研究也表明，同一物种分布范围内的可迁移程度仍可随环境背景变化\cite{lovell2025}。此外，清单报告预测也不是对潜在占域模型的独立验证。生态解释与报告预测分别为监测决策提供不同信息\cite{elith2009}。在鸟类预测研究中，Ovaskainen 等\cite{ovaskainen2026}将公众科学录音与标准化间隔记录、固定定点调查网络结合，展示了可重复观测设计对预测应用的支持。

本研究考察不同海拔与绿度环境中的报告回升程度，检验这些环境是否也对应较高的占域持续概率，并据此讨论环境差异的监测方法。首先，结合缺失组成与标准化报告预测，区分环境代表性和层内时期变化。其次，校正动态占域模型的时期标签对齐，将同一环境组合的报告轨迹与占域持续概率、定殖和重复报告相对照。最后，利用空间、时间验证及等预算分层排序，判断报告预测能否为需要复访的环境补充线索。常规模型比较和探索性神经流程检查均服务于这一监测问题。本研究使用同一公开观测集进行二次分析，其解释增量在于连接报告回升的环境分布、校正后的占域持续性关联与监测覆盖的实际取舍。
